\documentclass[a4paper,11pt]{article}

\usepackage[a4paper, top=2.5cm, bottom=2.5cm]{geometry}
\usepackage[utf8]{inputenc}
\usepackage[french, english]{babel}

\selectlanguage{french}

\usepackage{enumitem}
\usepackage{hyperref}
\usepackage{multicol}
\usepackage{float}
\usepackage{array}
\usepackage{caption}
\usepackage{listings}

%BEGIN LATEX
  \usepackage{core}

  \newcommand{\mousevox}[2]{\IfStrEq{\jobname}{\detokenize{mouse}}{#1}{#2}}

  % enumitem
  \setlength{\itemsep}{0\baselineskip}
  \setlength{\parsep}{0pt}
%END LATEX

\title{
  \blogtitle{Du Café Façon Guinness}
}
\author{}
\date{}

\pagestyle{plain}
\begin{document}
\neobar{article00.fr}{nav/navi.fr}

\maketitle

\thispagestyle{first}
\begin{blogmain}
\begin{multicols}{2}
  S'il y est un or noir, un carburant essentiel à tout dev, c'est le café.

  Cette graine aux propriétés psychotropes est ce qui nous permet de nous donner au maximum.

  Celle-ci est torréfiée depuis des siècles, probablement cinq tout au moins.

  Les méthodes de préparations divergent, il en existe aujourd'hui plus d'une dizaine.

  De quoi un dev a-t-il idéalement besoin ?
  \begin{itemize}
    \item d'un café prêt pour la semaine
    \item qui puisse être rapidement servi à tout moment
    \item qui soit sans sucre
    \item qui est un goût agréable, doux
  \end{itemize}

  \mousevox{
    Si nos critères n'étaient que cela, je conseillerais des pilules de café LP (à Libération Prolongée).

    Ajoutons donc des critères de confort :
    \begin{itemize}
      \item qui est un goût agréable, doux
      \item qui se rapproche d'un Irish-Coffee ou d'une Stout
    \end{itemize}
  }{}

  La plus récente est le "café à l'azote infusé à froid", celle-ci aurait été expérimentée vers 2010.
  Elle est l'évolution du café infusé à froid de Kyoto, et des progrès de l'industrie des boissons gazeuses.

  La méthode d'extraction est donc d'infuser de la moulure de café, généralement de 12 à 24h, à température de 0 à 8°C, puis de le filtrer, par exemple avec un erlenmeyer comme une Chemex qui peut permettre de filtrer jusqu'à 10 par 10 tasses

  \doubleimage{.6}{jar-monochrome}{Jar de moulure de café}{chemex-monochrome}{Chemex}

  Le café filtré sera ensuite versé dans un fût, celui-ci sera pressurisé avec de l'azote à l'état gazeux.
  Le café sera servi à la tireuse

  \simpleimage{.5}{gaz-monochrome}{Fût + Adapter/Régulateur + Bouteille}

  Son goût est crémeux, très comparable à celui d'une Guinness.
  Je suis un amateur de Guinness, j'ai découvert que dans les années 50, cette société a été la première à commercialiser de la bière pressurisée à l'azote.
  Un partenariat entre Guinness et Tiki Tonga Coffee nous permet aujourd'hui d'avoir du café aux arômes de Guinness.
  J'ai tellement aimé cela que j'utilise aujourd'hui une tasse graduée pour limiter ma consommation...

  % Café gradué

  L'azote qui permet de servir pressions ces cafés est distribué en cartouches ou en bouteilles.
  J'utilise personnellement une bouteille originellement destinée à de la Stout.

  Les bouteilles ont l'avantage d'être plus économique.
  Elles ont néanmoins des contraintes de stockages, à une température idéalement de 10°C à 30°C, strictement inférieur à 50°C, positionné verticalement, idéalement chaîné pour éviter qu'elle ne tombe.
  Le niveau de danger est un celui d'un gaz sous pression, non inflammable et non toxique.

  Lorsque j'ai commencé ce projet, j'étais ignorant de tout, je n'avais jamais moulu de café, j'ai reçu conseil d'un ami barman sur la bouteille d'azote, et d'autres sur la Chemex.
  Celui-ci me semble accessible pour toute personne curieuse qui aime le café.
\end{multicols}

Voici l'inventaire du matériel que j'utilise :
\begin{table}[htb]
  \centering
  \begin{tabular}{|l|c|c|}
    \hline \bf Description& \bf Fournisseur& \bf Réf-Fournisseur \\ \hline

    Cafetière Chemex en verre 10 tasses &
    \href{https://www.maxicoffee.com/chemex-m-197.html}{MaxiCoffee, Chemex} &
    \href{https://www.maxicoffee.com/cafetiere-chemex-verre-tasses-p-30943.html}{30943} \\

    Système mini fût 2L &
    \href{https://ikegger.eu/fr}{iKegger} &
    \href{https://ikegger.eu/fr/products/ikegger-2-0-complete-kickstarter-packages}{ich-zapfe} \\

    Adaptateur/Régulateur à double jauge &
    \href{https://ikegger.eu/fr}{iKegger} &
    \href{https://ikegger.eu/fr/products/nitrogen-bottle-adaptor-for-m3-dual-gauge-multi-gas-regulator}{adaptor-for-m3} \\

    Bouteille de gaz Azote (N2), 2L &
    \href{https://www.tireusesabiere.fr}{Tireuse à bière} &
    \href{https://www.tireusesabiere.fr/bouteille-d-azote-2-l-n2-avec-porte-bouteille/a-452064}{443527} \\

    (bonus) Siphon à Chantilly &
    \href{https://www.maxicoffee.com/siphon-isi-m-13.html}{MaxiCoffee, iSi} &
    \href{https://www.maxicoffee.com/siphon-chantilly-cream-profi-whip-inox-brosse-p-1333.html}{1333} \\

    \hline
  \end{tabular}
\end{table}

J'y ajoute en bonus de la crème Chantilly ou fouettée

Un projet suivant, sûrement bien plus facile, pourrait être de remplir des bouteilles d'eau gazeuses ou de soda avec du Co2.
\end{blogmain}
\end{document}
